\section*{Chapter \ref{ch:m3:physical-commodity-markets} --- Physical Commodity Markets}

\begin{solution}{exo:m3:physical-commodity-markets:1}
The five assessments sum to 403.25, an average of 80.65; plus 0.35, the price is
\qty{81.00}{\$/\barrel}.
\end{solution}

\begin{solution}{exo:m3:physical-commodity-markets:2}
$78.40 + 1.35 + 0.04 = \qty{79.79}{\$/\barrel}$ CIF. Under both terms the risk
passes to the buyer when the oil is on board at the loading port: the buyer bears
the voyage risk either way; under CIF the seller pays for the freight and the
insurance that covers the buyer.
\end{solution}

\begin{solution}{exo:m3:physical-commodity-markets:3}
$84.10 \times 0.998 - 2.20 - 0.05 = \qty{81.68}{\$/\barrel}$.
\end{solution}

\begin{solution}{exo:m3:physical-commodity-markets:4}
It buys 700 futures today and sells 140 on each pricing day. If a day is dropped the
price is the average of four assessments: it sells 175 on each of the four days.
\end{solution}

\begin{solution}{exo:m3:physical-commodity-markets:5}
\qty{1.08}{\$/\barrel}: the hedged cost is $F_0$ plus the average basis plus the
differential (\cref{prop:m3:physical-commodity-markets:hedged}), whose standard
deviation is proportional to the basis volatility. The futures moves are bought and
sold back on the same days as the physical prices in, so they cancel exactly.
\end{solution}

\begin{solution}{exo:m3:physical-commodity-markets:6}
The \omterm{def:m3:physical-commodity-markets:netback}{netbacks} are 80.17 (Rotterdam) and 80.02 (Singapore): Singapore's delivered
price must rise by more than \qty{0.15}{\$/\barrel}, to above
\qty{83.25}{\$/\barrel}.
\end{solution}

\begin{solution}{exo:m3:physical-commodity-markets:7}
September 2011, $-\qty{27.31}{\$/\barrel}$; WTI averaged above Brent in 4 of the
months from January 2011 to August 2026 (115 of all 320 months since 2000, almost all
before 2011).
\end{solution}

\begin{solution}{exo:m3:physical-commodity-markets:8}
The cargo's cost starts to float when the contract is signed, not on the B/L date,
and it fixes one fifth a day over the five pricing days, two of which come before
the B/L date. Buy the 700 futures when the exposure is created and sell 140 on each
pricing day at that day's settlement; selling them all ``once the price is known''
leaves the last days' futures moves unmatched by any physical move.
\end{solution}

\begin{solution}{pb:m3:physical-commodity-markets:1}
\textbf{1.} Zero: both prices float with Dated. \textbf{2.} It rises by 140\,000
barrels a day as the purchase prices in, to 700\,000 long on day 12. \textbf{3.}
700\,000 barrels long, from day 12 to day 27: 16 days. \textbf{4.} It falls by
140\,000 a day to zero on day 32. \textbf{5.} A trapezoid
(\cref{fig:m3:physical-commodity-markets:exposure}).
\textbf{6.} Sell 140 futures a day. \textbf{7.} Buy back 140 a day.
\textbf{8.} Short 700. \textbf{9.} A floating purchase against a floating sale has
no flat-price exposure. \textbf{10.} A swap on the Dated average of each \omterm{def:m3:physical-commodity-markets:formula}{pricing
period} (a Dated-to-futures swap, \cref{ch:m3:crude-oil}), or futures plus the
residual basis.
\textbf{11.} $1.10 - 0.20 - 0.65 - 0.03 = \qty{0.22}{\$/\barrel}$. \textbf{12.}
\$154\,000. \textbf{13.} Twenty business days between the \omterm{def:m3:physical-commodity-markets:formula}{pricing periods}:
$1.80\sqrt{20} \times 700\,000 \approx$ USD~5.63 million. \textbf{14.} The basis
between Dated and the futures at each pricing day, volume tolerance, demurrage and
delays, counterparty credit, and operational losses. \textbf{15.} The sale prices on
days 33--37: the exposure lasts 21 days at its maximum and the buy-backs move five
days later; the hedge must be kept open, and the futures may need rolling.
\textbf{16.} To capture the freight and the destination premium, and to control the
voyage. \textbf{17.} The purchase price, paid about thirty days after loading,
until the sale proceeds arrive, plus the margin on the futures. \textbf{18.} Its P\&L
is the differential between two prices, not the price level. \textbf{19.}
\emph{Named result:} \$154\,000 locked per cargo (\qty{0.22}{\$/\barrel}) against an
unhedged standard deviation of about USD~5.63 million. \textbf{20.} The spread
between a commodity's value in two places, times or forms, less the cost of moving
it, not the price level.
\end{solution}

\begin{solution}{iq:m3:physical-commodity-markets:1}
FOB: the seller loads onto the buyer's ship and the buyer pays freight and
insurance. CIF: the seller pays freight and insurance to the destination. In both the
risk passes on board at the loading port, so the buyer bears the voyage risk.
\iqlookfor{the distinction between who pays and who bears the risk.}
\end{solution}

\begin{solution}{iq:m3:physical-commodity-markets:2}
An average is harder to move by one trade or one bad day, spreads the pricing over
the time the cargo is actually loaded or delivered, and lets both sides hedge day by
day in liquid paper.
\iqlookfor{manipulation resistance, matching the operational window, hedgeability.}
\end{solution}

\begin{solution}{iq:m3:physical-commodity-markets:3}
At the trade date, hedge the whole volume, because the exposure exists as soon as
the price will float; then close $1/n$ of the hedge on each of the $n$ pricing days
(the month's business days). A monthly average swap does this in one instrument.
\iqlookfor{exposure created at signing, unwound by pricing days, the swap as the natural hedge.}
\end{solution}

\begin{solution}{iq:m3:physical-commodity-markets:4}
The settlement is computed by rule from exchange trades and is reproducible; the
assessment is a judgement by an editorial team on bids, offers and trades in a
window, sometimes with few trades, and can be revised. Its methodology, window and
timestamp matter, and its history is licensed data.
\iqlookfor{reproducibility, methodology risk, licensing.}
\end{solution}

\begin{solution}{iq:m3:physical-commodity-markets:5}
On basis between the benchmark hedged and the cargo's formula, on timing mismatches
(delays, tolerances), on counterparty default, on operational losses and demurrage,
and on the financing of margin calls on its hedges while the physical gain is
unrealised.
\iqlookfor{basis, operations, credit and liquidity rather than flat price.}
\end{solution}

\begin{solution}{iq:m3:physical-commodity-markets:6}
Deal capture with each leg's formula, benchmark and pricing days; a calendar of
business days per benchmark; a nightly job that computes the fixed fraction of each
leg, aggregates exposure by benchmark and day, compares it with futures and swaps
held, and outputs the trades to flatten; alerts on missing assessments and on legs
whose dates move.
\iqlookfor{the leg-and-pricing-day data model, per-benchmark aggregation, reconciliation with hedges.}
\end{solution}
